\section{协向量的拉回}
\label{sec:mmaps-pullback}

切向量随映射向前推，从$p$处送到$F(p)$处。余切向量是对切向量的读数，它的换算方向恰好相反：如果已经知道怎样在$F(p)$处读数，就可以先把$p$处的切向量推过去，再在那里读数，这样得到的是$p$处的一个读数。本节把这个想法写成定义，并说明方向相反带来的好处。

\paragraph{逐点的拉回}
设$F:M\to N$光滑，$p\in M$，$\omega\in T^*_{F(p)}N$。对$\bm v\in T_pM$，$\mathrm dF_p(\bm v)$属于$T_{F(p)}N$，可以被$\omega$读取。

\begin{definition}[余切向量的拉回]
\label{def:mmaps-covector-pullback}
$\mathrm dF_p$的对偶映射
\[
(\mathrm dF_p)^*:T^*_{F(p)}N\to T^*_pM,\qquad
\bigl((\mathrm dF_p)^*\omega\bigr)(\bm v)=\omega\bigl(\mathrm dF_p(\bm v)\bigr)
\]
称为$F$在$p$处的余切映射，$(\mathrm dF_p)^*\omega$称为$\omega$的\term{拉回}（pullback）。
\end{definition}

右端对$\bm v$线性，所以$(\mathrm dF_p)^*\omega$确实是$p$处的余切向量；它对$\omega$也线性。在坐标中，设$\omega=\sum_kb_k\,\mathrm dy_k|_{F(p)}$，$\bm J=\bm J_{\hat F}(\varphi(p))$。把$(\mathrm dF_p)^*\omega$作用在$\partial_j|_p$上，由式\eqref{eq:mmaps-differential-basis}得到它的第$j$个系数$\sum_kb_k\,\partial\hat F_k/\partial x_j$。写成列向量，
\begin{equation}
\bm a=\bm J^\top\bm b .
\label{eq:mmaps-pullback-coordinates}
\end{equation}
切向量的分量用$\bm J$从$M$推到$N$，余切向量的系数用$\bm J^\top$从$N$拉回$M$，两个方向相反。

\paragraph{协向量场的拉回总有定义}
对$N$上的协向量场$\omega$，在每个$p\in M$处拉回$F(p)$处的值，得到$M$上的协向量场
\[
(F^*\omega)_p=(\mathrm dF_p)^*\bigl(\omega_{F(p)}\bigr),\qquad p\in M .
\]
这里每个$p$只需要一个值$\omega_{F(p)}$，它总是存在且唯一。所以向量场推送时遇到的两类困难都不出现：多个原像只是各自拉回同一个值，没有原像的点在$M$中根本不出现。在坐标中，$F^*\omega$的系数是$\bm J_{\hat F}(\bm x)^\top\hat{\bm b}(\hat F(\bm x))$，由光滑函数复合、相乘得到，所以$\omega$光滑时$F^*\omega$也光滑。拉回对任何光滑映射都适用，不需要$F$可逆。

\paragraph{运算规则}
拉回与函数、复合和微分的关系可以归结为三条规则。

\begin{theorem}[拉回的运算规则]
\label{thm:mmaps-pullback-rules}
设$F:M\to N$与$G:N\to P$光滑，$\omega$是$N$上的光滑协向量场，$g,f\in C^\infty(N)$。则
\begin{enumerate}
\item $F^*(g\,\omega)=(g\circ F)\,F^*\omega$；
\item $(G\circ F)^*=F^*\circ G^*$；
\item $F^*(\mathrm df)=\mathrm d(f\circ F)$。
\end{enumerate}
\end{theorem}

\begin{proof}
第一条：在$p$处，$g\,\omega$的值是$g(F(p))\,\omega_{F(p)}$，对偶映射线性，系数$g(F(p))$直接提出。第二条：由链式法则，$\mathrm d(G\circ F)_p=\mathrm dG_{F(p)}\circ\mathrm dF_p$；对偶映射交换复合次序，$(\mathrm d(G\circ F)_p)^*=(\mathrm dF_p)^*\circ(\mathrm dG_{F(p)})^*$。第三条：对$\bm v\in T_pM$，
\[
\bigl(F^*\mathrm df\bigr)_p(\bm v)=\mathrm df_{F(p)}\bigl(\mathrm dF_p(\bm v)\bigr)=\mathrm d(f\circ F)_p(\bm v),
\]
最后一步是链式法则应用于$f\circ F$。
\end{proof}

第三条说明“先微分再拉回”与“先复合再微分”给出同一个结果。复合次序在第二条中反过来：先作用$F$再作用$G$的映射，拉回时先用$G^*$再用$F^*$。这与切向量的推送$\mathrm d(G\circ F)=\mathrm dG\circ\mathrm dF$方向相反，是对偶映射的一般性质。

\paragraph{例子：圆周上的角度读数}
圆周上的两张角度图相差常数平移，所以$\mathrm d\theta_1$与$\mathrm d\theta_2$在重叠处相等，拼成整个圆周上的光滑协向量场，记为$\mathrm d\theta$。注意它不是某个全局光滑函数的微分，因为圆周上没有处处光滑的全局角度函数，记号只表示它在每张图中都是坐标函数的微分。

用$E(t)=(\cos t,\sin t)$拉回：$E$在角度图中的坐标表示是$t\mapsto t+c$，导数为$1$，由式\eqref{eq:mmaps-pullback-coordinates}，$E^*(\mathrm d\theta)=\mathrm dt$。这里$E$不是单射，每个圆周上的点有无穷多个原像，拉回照样良定义。这一计算也证实了上一段的断言：若$\mathrm d\theta=\mathrm df$对某个$f\in C^\infty(S^1)$成立，则由第三条规则$\mathrm d(f\circ E)=E^*\mathrm d\theta=\mathrm dt$，于是$f(E(t))=t+c$；但$f\circ E$以$2\pi$为周期，矛盾。与上一节对比，$\mathbb R$上的向量场$t\frac{\mathrm d}{\mathrm dt}$无法推送到圆周，而圆周上的协向量场总能拉回$\mathbb R$。

\paragraph{例子：旋转下的高度读数}
由运算规则的第三条，$R_\theta^*(\mathrm dh)=\mathrm d(h\circ R_\theta)=\mathrm dh$，因为旋转不改变高度。这与上一节在$p=(0,\tfrac35,\tfrac45)$处的矩阵计算是同一件事：那里用行向量左乘$\bm R(\theta)$，正是式\eqref{eq:mmaps-pullback-coordinates}的转置形式。用拉回的语言，计算不需要逐点核对，直接由函数的不变性得出。

\paragraph{向更一般对象的推广}
本节只拉回了余切向量，即对单个切向量的线性读数。同样的思路可以用于同时读取多个切向量的多重线性读数，例如面积元和体积元，其中最重要的是交替的情形，即微分形式。它们的拉回、外微分及与积分的关系，在第~\ref{chap:manifold-calculus}章定义交替形式之后建立；黎曼度量作为对两个切向量的对称读数，其拉回在第~\ref{chap:riemannian-geometry}章用于定义等距映射。
